\section{Apresentação e Discussão dos Resultados: Plano de Análise}

Nesta etapa de qualificação, as regressões finais ainda não foram estimadas. A seção estabelece como seus resultados serão apresentados e interpretados após a construção e a validação do painel.

A discussão começará pela cobertura efetiva da amostra: municípios e anos observados, composição por bioma, presença da soja e perdas decorrentes de ausências ou incompatibilidades territoriais. Serão apresentadas distribuições de vegetação e clima, de modo a mostrar onde as comparações encontram suporte. Médias estaduais não substituirão a análise da sobreposição entre grupos, sobretudo quando a pequena amostra pantaneira limitar H2-a.

Na sequência, serão apresentados os efeitos marginais de déficit e excesso de chuva em diferentes estoques vegetais, acompanhados de intervalos de confiança e magnitudes em unidades produtivas. A discussão de H2 considerará tanto o teste conjunto das interações quanto o tamanho dos contrastes. Para H2-a, serão comparadas formalmente as respostas entre grupos, evitando confundir diferenças de precisão com diferenças dos efeitos.

Se maior vegetação estiver associada a menor perda produtiva diante de um choque adverso, será discutida sua compatibilidade com os mecanismos de água, solo e paisagem examinados na literatura. Se a associação indicar intensificação da perda, serão consideradas diferenças de localização, manejo, competição por recursos e mensuração da cobertura. Se os efeitos forem pequenos ou imprecisos, serão distinguidas a ausência de magnitude economicamente relevante e a insuficiência de informação para estimá-la. Em todos os casos, a discussão dependerá das condições de identificação sustentadas pelo desenho.

As implicações para conservação serão proporcionais ao que tiver sido identificado. Diferenças de resposta produtiva podem indicar que a condição ambiental e as particularidades territoriais merecem integrar a avaliação dos incentivos enfrentados pelos produtores. Elas não bastam, contudo, para fixar pagamentos por serviços ambientais, preços de carbono, penalidades ou metas ótimas de conservação. Essas escolhas também exigem custos de oportunidade, efeitos externos, distribuição dos benefícios e condições institucionais que não são integralmente observadas neste artigo.

Por fim, a relação com o capítulo 1 será retomada como articulação da tese: os incentivos à conversão podem afetar um estoque que participa das condições futuras de produção. Essa interpretação preserva a dimensão intertemporal do problema, mas não atribui às estimativas de H2 um teste direto de expectativas, taxas de desconto ou decisões dinâmicas de conservação.
